\documentclass{article}

% ----------------------- PAQUETES ---------------------- %
% Set the font (output) encodings
\usepackage[T1]{fontenc}
\usepackage[spanish]{babel}
\usepackage{graphicx}							% paquete para poder usra graficos
\usepackage{float}								% paquete para poder usar \begin {figure}[H]
\usepackage{hyperref}							% paquete para hacer hyperreferencias a links
\usepackage{fancyhdr}							% paquete para poner cosas en footer y header
\usepackage{wrapfig}							% paquete para hacer wrapfigure
\usepackage{titling}							% paquete para modificar el espaciado de los titulos
\usepackage{titlesec}							% paquete para modificar como se ven los titulos
\usepackage{minted}								% paquete para agregar codigo al documento
\usepackage[smartEllipses]{markdown}							% paquete para poder usar markdown para incluir las bitacoras


% ------------------------------------------------------- %

% ---------------------- CONFIGURACION ------------------ %

\setlength{\headheight}{12pt} %esto era para evitar un error de hbox
\graphicspath{{images/}}

% ----------------- Configuracion de hyperref ----------- %
\hypersetup{								
	colorlinks=true,
	linkcolor=black,			%modo claro
	%linkcolor=white,		%modo oscuro
	filecolor=brown,		
	urlcolor=blue,
	pdftitle={Proyecto Indoor},
	}
% ------------------------------------------------------- %

% ----------------- Configuracion de fancyhdr --------------- %
\pagestyle{fancy}
\setlength{\headheight}{12.17pt} %esto esta para solucionar un warning del fancyhdr
\rfoot{Página \thepage}
\cfoot{}
\lfoot{Krapp Ramiro}
\renewcommand{\sectionmark}[1]{%
\markboth{\thesection\quad #1}{}}
\fancyhead[L]{\leftmark}
\fancyhead[R]{Traintorio}
\renewcommand{\headrulewidth}{2pt}
\renewcommand{\footrulewidth}{1pt}
% ------------------------------------------------------- %


%-------------- Formatos de los titulos --------------%
\newcommand{\sectionbreak}{\clearpage}

\titleformat{\section}
	{\bfseries \huge}
	{}
	{0em}
	{}[\titlerule]

\titleformat{\subsection}
	{\bfseries \Large}
	{}
	{0em}
	{}

\titleformat{\subsubsection}
	{\bfseries \large}
	{}
	{0em}
	{}

\titlespacing{\section} %me permite controlar el espaciado de la seccion que le indico
	{0em}
	{0em}
	{1.5em}

\titlespacing{\subsection}
{0em} %sangria
{3em} %separacion con lo que hay arriba
{0.5em} %separacion con lo que hay abajo
%----------------------------------------------------- %

% ----------------- FIN DE CONFIGURACION ------------- %


% Spanish-specific commands
\begin{document}
\begin{titlepage}
	\begin{center}
		\vspace{1cm}

		{\Huge
			\textbf{Traintorio}}

		\vspace{0.3cm}
		{\LARGE
			Informe técnico}

		\vspace{0.5cm}
		{\Large
			Un trabajo presentado para la materia de \\
			Proyectos y Diseño Electrónico}

		\vspace{2cm}

		\begin{figure}[H]
			\centering
			\includegraphics[width=0.6\textwidth]{logo.png}
		\end{figure}

		\vfill

		{\Large
			\textbf{Krapp Ramiro} \\
			\vspace{0.5cm}
			Instituto tecnológico San Bonifacio\\
			Departamento de electrónica\\
			\today
		}

		\vspace{0.5cm}
		{\large Hecho en {\LaTeX}\\
			Versión Alpha 0.1}

	\end{center}
\end{titlepage}

\tableofcontents		\noindent\rule{\textwidth}{0.7pt}
El índice tiene hipervínculos incorporados!
Toca en cada seccion y automaticamente tu lector de pdfs te llevara a esa página
\\[12pt] %esto esta para generar un espaciado entre lineas
{\large Tengo un
\href{https://github.com/KrappRamiro/traintorio}{Repositorio en GitHub}}\\
{\small \url{https://github.com/KrappRamiro/traintorio}}



\section{Introducción}

\section{Diagrama esquematico}
%\includegraphics[width=0.9\textwidth]{../../diagrama_drawio.png}

\section{Base de datos}
%\includegraphics[width=0.9\textwidth]{../../database_schema.png}

\section{Codigo del programa}
%\inputminted %TODO: ponerle inconsolata
%[frame= lines, linenos, breaklines, tabsize = 3, fontsize=\footnotesize,
%	label=Codigo principal]
%{cpp}{../src/main.cpp}

\section{Bitacoras Personales}

% \markdownInput{../bitacoras/bitacora_krapp.md}
\end{document}